\section{Задание}

Используя хвостовую рекурсию, разработать программу, позволяющую найти
\begin{itemize}
    \item n!
    \item n-e число Фибоначчи 
\end{itemize}

Убедиться в правильности результатов.
Для одного из вариантов вопроса и каждого задания составить таблицу, отражающую
конкретный порядок работы системы.
Т.к. резольвента хранится в виде стека, то состояние резольвенты требуется отображать
в столбик: вершина – сверху! Новый шаг надо начинать с нового состояния резольвенты!
Для одного из вариантов вопроса составить таблицу, отражающую конкретный порядок
работы системы.
Для одного из вариантов вопроса и конкретной БЗ составить таблицу, отражающую кон-
кретный порядок работы системы, с объяснениями: очередная проблема на каждом шаге и
метод ее решения; каково новое текущее состояние резольвенты, как получено; какие дальней-
шие действия? (Запускается ли алгоритм унификации? Каких термов? Почему этих?) ; вывод
по результатам очередного шага и дальнейшие действия.